\documentclass[11pt,a4paper]{article}

\usepackage[utf8]{inputenc}
\usepackage[T1]{fontenc}
\usepackage{lmodern}
\IfFileExists{brazil.ldf}{\usepackage[brazil]{babel}}{%
  \renewcommand{\figurename}{Figura}
  \renewcommand{\tablename}{Tabela}
  \renewcommand{\contentsname}{Sumário}
}
\usepackage[margin=2.5cm]{geometry}
\usepackage{microtype}
\usepackage{graphicx}
\usepackage{booktabs}
\usepackage{tabularx}
\usepackage{array}
\usepackage{float}
\usepackage[font=small,labelfont=bf]{caption}
\usepackage{xcolor}
\usepackage{listings}
\usepackage[hidelinks]{hyperref}

\lstset{basicstyle=\ttfamily\small, columns=fullflexible, frame=single, breaklines=true}
\newcolumntype{Y}{>{\raggedright\arraybackslash}X}

\newcommand{\figura}[3][\textwidth]{%
  \begin{figure}[H]
    \centering
    \IfFileExists{figuras/#2.png}{%
      \includegraphics[width=#1,height=0.75\textheight,keepaspectratio]{figuras/#2.png}}{%
      \fbox{\parbox[c][4cm][c]{0.9\textwidth}{\centering\ttfamily figuras/#2.png}}}
    \caption{#3}
    \label{fig:#2}
  \end{figure}}

\title{Processador RISC-V RV32I de ciclo único\\[0.3em]
  \large Relatório final}
\author{Gabriel Sena Barbosa Silva\\
  Jair Rocha Souza\\
  João Vítor Cedraz Carneiro\\
  Luiz Augusto dos Santos Oliveira\\
  Tássio Carvalho Rodrigues\\[0.8em]
  Programa CI Digital --- CEPEDI}
\date{31 de outubro de 2026}

\begin{document}
\maketitle

\noindent Repositório: \texttt{Gabriel-SBS/rv32i-processor}

\tableofcontents

\section{Arquitetura}

O processador implementa o conjunto de instruções base RV32I, de 32 bits, em ciclo único:
cada instrução é buscada, decodificada, executada, acessa a memória quando precisa e grava
o resultado no banco de registradores dentro de um único ciclo de clock.

Para explicar o caminho de dados, dividimos o processador nos cinco estágios da
Figura~\ref{fig:arq_visao_geral}. Não há registradores entre eles, porque o processador é de
ciclo único; a divisão é só para organizar a explicação. Nas figuras de cada estágio, os blocos com contorno contínuo
pertencem ao estágio, e os de contorno tracejado são sinais que vêm de outros estágios ou
vão para eles. Os códigos Mermaid estão em \texttt{relatorio/figuras/arq\_*.mmd}.

\figura{arq_visao_geral}{Visão geral: os cinco estágios e o que passa entre eles. A linha
pontilhada é a volta para a busca, usada nos desvios e no JALR.}

\subsection{Clock e reset}

O reset \texttt{rst\_n} é ativo em nível baixo e assíncrono. Com \texttt{rst\_n = 0}, o PC vai
para \texttt{0x00000000} na hora, sem esperar o clock. As escritas acontecem na borda de
subida de \texttt{clk}. O registrador \texttt{x0} vale sempre zero: escritas nele são ignoradas.

\subsection{Busca}

\figura{arq_busca}{Busca da instrução e cálculo do próximo PC.}

O PC endereça a memória de instruções, que tem 256 palavras de 32 bits, lidas de
\path{instructions/instruction_control.hex} no início da simulação. O endereço é
\texttt{pc[9:2]}, então os bits 1:0 são ignorados. No \texttt{tb\_processor}, o programa
escolhido com \texttt{+TEST=<nome>} é carregado por cima desse arquivo.

Ao mesmo tempo, o \texttt{pc\_logic} calcula os candidatos a próximo PC, e um \texttt{mux4}
escolhe entre eles pelo sinal \texttt{PC\_sel}:

\begin{itemize}
  \item \texttt{00}: \texttt{PC + 4}, na execução normal e quando o desvio não é tomado;
  \item \texttt{01}: \verb|(rs1 + imm) & ~1|, no \texttt{JALR}, com o resultado vindo da ULA;
  \item \texttt{10}: \texttt{PC + imm}, nos desvios tomados e no \texttt{JAL}.
\end{itemize}

A seleção \texttt{11} não é usada e também dá \texttt{PC + 4}. O valor \texttt{PC + 4} vai
ainda para o write-back, porque é o endereço de retorno que o JAL e o JALR gravam em
\texttt{rd}.

\subsection{Decodificação}

\figura{arq_decodificacao}{Decodificação. Cada campo da instrução vai para um bloco
diferente.}

A instrução é quebrada em campos, e cada campo vai para um bloco:

\begin{itemize}
  \item \texttt{controller} lê o \texttt{opcode} e gera \texttt{we\_reg}, \texttt{we\_mem},
        \texttt{rs1\_sel}, \texttt{rs2\_sel}, \texttt{REG\_sel}, \texttt{PC\_sel} e
        \texttt{ALU\_Op}. Nos desvios, quem diz se o salto acontece é a saída
        \texttt{branch\_taken} da ULA: com ela em 1, o \texttt{controller} põe
        \texttt{PC\_sel = 10}.
  \item \texttt{alu\_controller} traduz \texttt{ALU\_Op}, \texttt{funct3} e \texttt{funct7[5]}
        para o código de operação da ULA.
  \item O banco de registradores tem 32 registradores de 32 bits, duas portas de leitura
        combinacional (\texttt{rs1} e \texttt{rs2}) e uma de escrita (\texttt{rd}), que só
        grava na borda de subida e com \texttt{we\_reg = 1}.
  \item \texttt{imm\_gen} monta o imediato dos formatos I, S, B, U e J, com extensão de sinal.
\end{itemize}

\subsection{Execução}

\figura{arq_execucao}{Escolha dos operandos e operação na ULA.}

Dois muxes escolhem as entradas da ULA. O \texttt{src\_a} é \texttt{rs1}, ou o PC quando
\texttt{rs1\_sel = 1} (AUIPC e JAL). O \texttt{src\_b} é \texttt{rs2}, ou o imediato quando
\texttt{rs2\_sel = 1} (instruções do tipo I, loads, stores, LUI, AUIPC, JAL e JALR).

A ULA é combinacional, com a interface \texttt{ALU (a, b, ALUControl, result, branch\_taken)}.
A codificação segue o próprio RISC-V: nas instruções do tipo R,
\texttt{ALUControl = \{funct7[5], funct3\}}. Os 16 códigos são usados
(Tabela~\ref{tab:ula}).

\begin{table}[H]
  \centering
  \caption{Códigos de operação da ULA.}
  \label{tab:ula}
  \begin{tabular}{cl@{\hspace{2.5em}}cl}
    \toprule
    \texttt{ALUControl} & Operação & \texttt{ALUControl} & Operação \\
    \midrule
    \texttt{0000} & ADD              & \texttt{1000} & SUB \\
    \texttt{0001} & SLL              & \texttt{1001} & BEQ (\texttt{a == b}) \\
    \texttt{0010} & SLT (com sinal)  & \texttt{1010} & BNE (\texttt{a != b}) \\
    \texttt{0011} & SLTU (sem sinal) & \texttt{1011} & BGE (com sinal) \\
    \texttt{0100} & XOR              & \texttt{1100} & BLT (com sinal) \\
    \texttt{0101} & SRL              & \texttt{1101} & SRA \\
    \texttt{0110} & OR               & \texttt{1110} & BLTU (sem sinal) \\
    \texttt{0111} & AND              & \texttt{1111} & BGEU (sem sinal) \\
    \bottomrule
  \end{tabular}
\end{table}

Os deslocamentos usam só \texttt{b[4:0]}. As comparações devolvem 1 ou 0 em \texttt{result}.
A saída \texttt{branch\_taken} só sobe nos seis códigos de desvio, e apenas quando a condição
é verdadeira. Num \texttt{SLT} que dá 1, por exemplo, ela continua em 0, então nenhuma
instrução que não seja desvio consegue mudar o fluxo do programa.

O \texttt{result} vai para a memória de dados, como endereço, e para o write-back. No JALR,
ele também é o destino do salto.

\subsection{Memória}

\figura[0.75\textwidth]{arq_memoria}{Memória de dados e os sinais que ela recebe.}

A memória de dados tem 512 palavras de 32 bits (2~KiB). O endereço, em bytes, vem do
\texttt{result} da ULA, e o dado a gravar vem de \texttt{rs2}. O \texttt{funct3} da
instrução diz o tamanho do acesso: byte (\texttt{sb}, \texttt{lb}, \texttt{lbu}), meia palavra
(\texttt{sh}, \texttt{lh}, \texttt{lhu}) ou palavra (\texttt{sw}, \texttt{lw}). A escrita é
síncrona, só com \texttt{we\_mem = 1}, e muda apenas os bytes do acesso. A leitura é
combinacional: \texttt{lb} e \texttt{lh} estendem o sinal, e \texttt{lbu} e \texttt{lhu}
completam com zeros.

\subsection{Write-back}

\figura[0.75\textwidth]{arq_writeback}{Mux de write-back, que escolhe o valor gravado em
\texttt{rd}.}

O último mux escolhe, pelo \texttt{REG\_sel}, o que será gravado em \texttt{rd}:

\begin{itemize}
  \item \texttt{00}: o dado lido da memória, nos loads;
  \item \texttt{01}: o \texttt{result} da ULA, nas instruções dos tipos R e I e no AUIPC;
  \item \texttt{10}: o imediato, no LUI;
  \item \texttt{11}: \texttt{PC + 4}, no JAL e no JALR.
\end{itemize}

Stores e desvios não gravam nada. Nesses casos o \texttt{controller} deixa
\texttt{we\_reg = 0} e a saída do mux é ignorada.

\section{Instruções implementadas}

\begin{table}[H]
  \centering
  \caption{Instruções do RV32I implementadas e testadas.}
  \begin{tabularx}{\textwidth}{lcY}
    \toprule
    Grupo & Formato & Instruções \\
    \midrule
    Aritméticas e lógicas entre registradores & R & ADD, SUB, AND, OR, XOR, SLL, SRL, SRA, SLT, SLTU \\
    Aritméticas e lógicas com imediato & I & ADDI, ANDI, ORI, XORI, SLLI, SRLI, SRAI, SLTI, SLTIU \\
    Loads & I & LB, LH, LW, LBU, LHU \\
    Stores & S & SB, SH, SW \\
    Desvios condicionais & B & BEQ, BNE, BLT, BGE, BLTU, BGEU \\
    Saltos & J / I & JAL, JALR \\
    Imediato superior & U & LUI, AUIPC \\
    \bottomrule
  \end{tabularx}
\end{table}

\section{Verificação}

A verificação teve duas partes.

Primeiro, cada bloco foi testado sozinho. São 14 testbenches, que combinam casos escritos à
mão (valores-limite, estouro, reset, \texttt{x0}) com milhares de entradas aleatórias
comparadas a um modelo de referência escrito dentro do próprio testbench.

Depois, o processador inteiro executou 10 programas em assembly RV32I. Cada programa termina
num laço \texttt{fim: j fim}. Quando o processador chega nesse laço, o \texttt{tb\_processor}
confere:

\begin{itemize}
  \item se \texttt{x31 = 1}, valor que o próprio programa grava quando todos os seus subtestes passam;
  \item se o PC parou no endereço esperado;
  \item se os 32 registradores têm os mesmos valores calculados pelo simulador de referência
        \texttt{rv32i\_sim.py}.
\end{itemize}

Os programas usam \texttt{x30} para guardar o número do subteste em andamento. Se algum falhar,
o programa grava \texttt{0xBAD} em \texttt{x31}, e \texttt{x30} mostra em qual subteste parou.

O código-fonte dos programas (\texttt{.s}), o montador \texttt{rv32i\_asm.py}, o simulador
\texttt{rv32i\_sim.py} e o script \texttt{gerar\_todos.sh}, que gera os \texttt{.hex} e os
\texttt{.regs}, estão em \path{instructions/arquivos_assembly/Treino_RISC_V/testes/}.

Todos os testes rodam com \texttt{sim/run\_xrun.sh} no Cadence Xcelium 25.03.

\section{Resultados}

\subsection{Testes de cada bloco}

\begin{table}[H]
  \centering
  \small
  \caption{Testbenches dos blocos. Todos passaram.}
  \begin{tabularx}{\textwidth}{llrY}
    \toprule
    Bloco & Testbench & Verif. & O que é testado \\
    \midrule
    Somador & \texttt{tb\_adder} & 8 & PC + 4, imediato positivo e negativo, estouro de 32 bits \\
    Mux 2:1 & \texttt{tb\_mux2} & 207 & as duas seleções, 200 casos aleatórios \\
    Mux 4:1 & \texttt{tb\_mux4} & 7 & as quatro seleções \\
    PC & \texttt{tb\_pc} & 11 & reset assíncrono, carga só na borda de subida \\
    Próximo PC & \texttt{tb\_pc\_logic} & 18 & sequencial, desvio tomado e não tomado, JAL, JALR \\
    ULA & \texttt{tb\_alu} & 51 & as 16 operações, os 6 desvios, \texttt{branch\_taken = 0} fora dos desvios \\
    ULA contra modelo & \texttt{tb\_alu\_ref} & 23.400 & 3.400 casos de borda e 20.000 aleatórios \\
    Banco de registradores & \texttt{tb\_register\_file} & 1.214 & duas portas de leitura, \texttt{x0}, \texttt{we = 0}, 500 aleatórios \\
    Banco + ULA & \texttt{tb\_regfile\_alu} & 5.168 & instruções que usam o resultado da anterior, 5.000 aleatórias \\
    Memória de dados & \texttt{tb\_data\_memory} & 3.559 & byte, meia palavra e palavra, 2.000 acessos aleatórios \\
    Memória de instruções & \texttt{tb\_instruction\_memory} & 46 & leitura alinhada, posições vazias, endereço além do fim \\
    Gerador de imediatos & \texttt{tb\_imm\_gen} & 4.526 & formatos I, S, B, U e J, 500 casos por opcode \\
    Controle da ULA & \texttt{tb\_alu\_controller} & 329 & todas as combinações de \texttt{ALU\_Op}, \texttt{funct3} e \texttt{funct7[5]} \\
    Unidade de controle & \texttt{tb\_controller} & 274 & 9 classes de instrução, os 128 opcodes \\
    \midrule
    Total & & 38.818 & \\
    \bottomrule
  \end{tabularx}
\end{table}

\subsection{Programas no processador completo}

\begin{table}[H]
  \centering
  \small
  \caption{Programas executados no processador. Todos passaram, com o mesmo número de ciclos
  do simulador de referência.}
  \begin{tabularx}{\textwidth}{lrrY}
    \toprule
    Programa & Ciclos & Verif. & O que exercita \\
    \midrule
    01 \texttt{alu\_reg} & 70 & 34 & operações entre registradores \\
    02 \texttt{alu\_imm} & 64 & 34 & operações com imediato \\
    03 \texttt{shifts} & 77 & 34 & deslocamentos lógicos e aritméticos \\
    04 \texttt{comparacoes} & 51 & 34 & comparações com e sem sinal \\
    05 \texttt{lui\_auipc} & 55 & 34 & LUI e AUIPC \\
    06 \texttt{load\_store\_word} & 84 & 34 & LW e SW com vários offsets \\
    07 \texttt{load\_store\_byte\_half} & 90 & 34 & LB, LH, LBU, LHU, SB, SH \\
    08 \texttt{branches} & 113 & 34 & os seis desvios condicionais \\
    09 \texttt{jal\_jalr} & 66 & 34 & JAL, JALR e endereço de retorno \\
    10 \texttt{algoritmo\_final} & 723 & 34 & soma de vetor, Fibonacci, bubble sort, contagem de bits e tamanho de string \\
    \midrule
    Total & 1.393 & 340 & \\
    \bottomrule
  \end{tabularx}
\end{table}

\subsection{Erros encontrados pelos programas}

Os blocos passavam nos testes individuais, mas três programas acharam erros que só aparecem
com tudo ligado:

\begin{table}[H]
  \centering
  \small
  \caption{Erros encontrados e correções.}
  \begin{tabularx}{\textwidth}{lYY}
    \toprule
    Programa & O que acontecia & Correção \\
    \midrule
    05 \texttt{lui\_auipc} & O AUIPC dava \texttt{PC + 2$\cdot$imm}: o mux de \texttt{src\_a}
      recebia \texttt{pc\_target}, que já é \texttt{PC + imm}, em vez de \texttt{pc}
      & \texttt{src\_a} passou a receber \texttt{pc} quando \texttt{rs1\_sel = 1} \\
    07 \texttt{load\_store\_byte\_half} & LB, LH, LBU, LHU, SB e SH não funcionavam: a memória só
      lia e escrevia palavras inteiras
      & a memória passou a receber \texttt{funct3} e o endereço em bytes \\
    10 \texttt{algoritmo\_final} & Com 64 palavras, os endereços \texttt{0x400} e \texttt{0x500}
      caíam na mesma posição
      & memória de 512 palavras, parâmetro \texttt{DM\_ADDR\_W = 9} \\
    \bottomrule
  \end{tabularx}
\end{table}

\section{Formas de onda}

As figuras foram geradas com o script \texttt{sim/figuras\_xrun.sh}. Ele simula cada programa no
\texttt{xrun} e para a simulação com um breakpoint condicional na instrução de interesse
(por exemplo, ``opcode igual a AUIPC''). A parada acontece na borda de descida do clock, no meio
do ciclo, quando todos os sinais já estão estáveis.

\begin{lstlisting}[language=bash]
sim/figuras_xrun.sh              # simula todas as figuras
sim/figuras_xrun.sh view fig04   # abre a fig04 com o cursor no ponto
\end{lstlisting}

Os sinais aparecem em grupos: clock e reset, busca, decodificação, controle, execução,
memória e status do teste (\texttt{x30} e \texttt{x31}). O clock tem período de 10~ns e o
reset é liberado em 20~ns.

\figura{fig01}{Fim do reset (programa 01). Em 20~ns \texttt{rst\_n} sobe com o PC em
\texttt{0x00000000}. A primeira instrução, \texttt{0x00000f93} (\texttt{addi x31, x0, 0}),
já está na saída da memória de instruções, e o PC passa a avançar de 4 em 4.}

\figura{fig02}{SUB (programa 01, 260~ns). A instrução \texttt{0x402081b3} é
\texttt{sub x3, x1, x2}, no PC \texttt{0x60}, com \texttt{x1 = 0xA} e \texttt{x2 = 3}.
Como \texttt{funct7 = 0x20}, o \texttt{alu\_controller} gera o código de SUB (\texttt{8})
e a ULA devolve 7.
\texttt{REG\_sel = 01} manda o resultado para \texttt{x3}.}

\figura{fig03}{ADDI com imediato negativo (programa 02, 80~ns). \texttt{0x80000093} é
\texttt{addi x1, x0, -2048}. O \texttt{imm\_gen} estende o sinal e \texttt{imm\_ext} vale
\texttt{0xfffff800}. O campo \texttt{funct7} aparece como \texttt{0x40}, mas no tipo I esses
bits fazem parte do imediato: \texttt{ALU\_Control} fica em 0 (ADD) e o resultado é
\texttt{0xfffff800}.}

\figura{fig04}{AUIPC (programa 05, 310~ns). \texttt{0x00000097} é \texttt{auipc x1, 0}, no PC
\texttt{0x78}. Com \texttt{rs1\_sel = 1}, a entrada \texttt{src\_a} da ULA recebe o PC e o
resultado é \texttt{0x78}. Era exatamente aqui que aparecia o erro do \texttt{PC + 2$\cdot$imm}.}

\figura{fig05}{SB (programa 07, 520~ns). \texttt{0x002502a3} é \texttt{sb x2, 5(x10)}, com
\texttt{x10 = 0x400} e \texttt{x2 = 0xAA}. A ULA calcula o endereço \texttt{0x405},
\texttt{we\_mem} vai a 1 e, na borda de subida, só o byte desse endereço é escrito.}

\figura{fig06}{LB (programa 07, 80~ns). \texttt{0x00050103} é \texttt{lb x2, 0(x10)}, com
endereço \texttt{0x400}. O byte guardado ali é \texttt{0x01}, então \texttt{mem\_o} vale
\texttt{0x00000001}, e \texttt{REG\_sel = 00} grava esse valor em \texttt{x2}.}

\figura{fig07}{Desvio tomado (programa 08, 90~ns). \texttt{0x00208463} é
\texttt{beq x1, x2, 8}, no PC \texttt{0x1c}. Os dois registradores valem 5, a ULA (código
\texttt{9}, BEQ) põe
\texttt{branch\_taken = 1} e o \texttt{controller} coloca \texttt{PC\_sel = 10}: o próximo PC
é \texttt{0x24}.}

\figura{fig08}{Desvio não tomado (programa 08, 110~ns). \texttt{0x10308663} é um
\texttt{beq x1, x3} no PC \texttt{0x28}, com \texttt{x1 = 5} e \texttt{x3 = 6}. Como os valores
são diferentes, \texttt{branch\_taken} fica em 0, \texttt{PC\_sel = 00} e a execução segue em
\texttt{0x2c}.}

\figura{fig09}{JAL (programa 09, 50~ns). \texttt{0x0d8000ef} é \texttt{jal x1, 0xd8}, no PC
\texttt{0x0c}. O próximo PC é \texttt{0xe4} (\texttt{PC + imm}) e \texttt{REG\_sel = 11}
grava em \texttt{x1} o endereço de retorno, \texttt{0x10}.}

\figura{fig10}{JALR (programa 09, 70~ns). \texttt{0x00008067} é \texttt{jalr x0, 0(x1)},
o retorno da sub-rotina, no PC \texttt{0xe8}. A ULA soma \texttt{x1 + 0}, o bit 0 é zerado e
a execução volta para \texttt{0x10} (\texttt{PC\_sel = 01}), logo depois do JAL da
Figura~\ref{fig:fig09}.}

\figura{fig11}{Fim do programa 10 (7246~ns). O PC chega ao laço final em \texttt{0x1a8}
(\texttt{0x0000006f}, \texttt{j fim}). \texttt{x30 = 5} mostra que os cinco subtestes rodaram
e \texttt{x31 = 1} indica que todos passaram.}

\section{Conclusão}

O processador RV32I de ciclo único está completo e executa todas as instruções da
Seção~2. Os 14 testbenches de bloco e os 10 programas passam no Xcelium, e em todos os
programas o estado final dos registradores e o número de ciclos batem com o simulador de
referência.

Os testes com programas completos foram os que mais ajudaram. Eles acharam três erros que os
testes de bloco não pegavam: no AUIPC, nos acessos de byte e meia palavra e no tamanho da
memória de dados. Cada bloco, sozinho, funcionava; os erros só apareceram com tudo ligado.

Tudo o que usamos está no repositório. O RTL e os testbenches ficam em \texttt{rtl/} e
\texttt{tb/}, e os programas em \texttt{instructions/programs/}, com os registradores
esperados em \texttt{.regs}. O código assembly e as ferramentas que geraram os testes estão em
\texttt{instructions/arquivos\_assembly/}, e os scripts de simulação e de formas de onda,
em \texttt{sim/}.

\end{document}
